% !PRINTED 584
\weberpage{584}

%<<CONT
zugrunde, die wir als Funktion von $u = v/2K$ mit $\boldsymbol{s}(u)$ bezeichnen. Diese Funktion hat dann die durch
\weq{4}{\begin{aligned}
\boldsymbol{s}(u+1) &= -\boldsymbol{s}(u), \\
\boldsymbol{s}(u+\omega) &= \boldsymbol{s}(u)
\end{aligned}}
ausgedrückte doppelte Periodizität. Alle Perioden von $\boldsymbol{s}(u)$ sind in der Form $2m + n\omega$ enthalten, worin $m, n$ ganze rationale Zahlen sind, und gehen also durch Multiplikation mit $a$ in ganze Zahlen über.

Zwei Zahlen $u, u'$, die sich nur um eine Periode unterscheiden, hei\ss en kongruent nach dem Modul $2, \omega$:
\[
u' \equiv u \ (\mathrm{mod}\ 2, \omega).
\]
\gesperrt{Da die Funktion $\mathrm{sn}\,v$ einen Wert nur zweimal im Periodenparallelogramm annimmt, so wird nur dann}
\[
\boldsymbol{s}(u) = \boldsymbol{s}(u'),
\]
\gesperrt{wenn entweder}
\weq{5}{\begin{aligned}
u &\equiv u' \\
&\text{oder} \\
u &\equiv 1 - u'
\end{aligned}\quad (\mathrm{mod}\ 2, \omega)}
\gesperrt{ist, und $\boldsymbol{s}(u)$ verschwindet nur, wenn}
\[
u \equiv 0, 1 \ (\mathrm{mod}\ 2, \omega)
\]
\gesperrt{ist.}

Da jede gebrochene Zahl in $\Omega$ durch Multiplikation mit einer ganzen rationalen Zahl in eine ganze Zahl verwandelt werden kann, so ergibt sich aus \S\,156, (6) der Satz:

1. \gesperrt{Ist $\eta$ irgend eine gebrochene Zahl des Körpers $\Omega$, so ist $\boldsymbol{s}(\eta)$ eine algebraische Zahl, und wenn insbesondere $\eta$ so dargestellt werden kann, da\ss{} sein Nenner relativ prim zu 2 ist, so ist $\boldsymbol{s}(\eta)$ eine ganze algebraische Zahl.}

Es sei jetzt
\weq{6}{\mu = y + x\sqrt{\varDelta}}
eine ganze Zahl in $\Omega$, in der $x, y$ ganze rationale Zahlen sind\footnote{Im Fall eines geraden $\varDelta$ kann hiernach $\mu$ jede ganze Zahl aus $\Omega$ sein. Im Fall eines ungeraden $\varDelta$ ist $\mu$ eine Zahl der Ordnung [2].}, und
\[
\mu' = y - x\sqrt{\varDelta},
\]
die zu $\mu$ konjugierte Zahl. Ferner
\weq{7}{m = \mu\mu' = y^2 - \varDelta x^2}
die Norm von $\mu$, die wir als \gesperrt{ungerade} voraussetzen.
